\documentclass[11pt]{article}
\usepackage{fontspec}
\usepackage[french,english]{babel}
\usepackage{amsmath}
\usepackage{amssymb}
\usepackage{amsthm}
\usepackage{mathrsfs}
\usepackage{tikz-cd}
\usepackage{enumitem}
\usepackage{lmodern}

\newcommand{\printedpage}[1]{\marginpar{\footnotesize p.~#1}}
\newcommand{\msp}[1]{\marginpar{\footnotesize [#1]}}

% short macros for printed notation
\newcommand{\uGamma}{\underline{\Gamma}}
\newcommand{\uFl}{\underline{\mathrm{Fl}}}
\newcommand{\uIs}{\underline{\mathrm{Is}}}
\newcommand{\uHom}{\underline{\mathrm{Hom}}}
\newcommand{\Ob}{\mathrm{Ob}\,}
\newcommand{\Fl}{\mathrm{Fl}\,}
\newcommand{\id}{\mathrm{id}}
\newcommand{\pr}{\mathrm{pr}}
\newcommand{\diag}{\mathrm{diag}\,}
\newcommand{\cardred}{\operatorname{card\,red}}
\newcommand{\cocart}{\mathrm{cocart}}
\newcommand{\cart}{\mathrm{cart}}
\newcommand{\cc}{\mathrm{c}}
\newcommand{\circled}[1]{\tikz[baseline=(c.base)]\node[draw,circle,inner sep=0.8pt](c){\upshape #1};}

\begin{document}

\printedpage{149}
% head: A. Grothendieck : Dérivateurs, Ch. XVIII, Édition des Universités de Montpellier II et Paris VII

\msp{page 393}
All these assertions are easily proved in the order in which they are given, using 4.14.17 and the general results already proved. Of course, we can begin by replacing $\mathcal{C}$ by an equivalent reduced subcategory $\mathcal{C}_0$ and $\mathcal{E}$ by $\mathcal{E}_0 = \mathcal{E}|\mathcal{C}$ \texttt{[rather} $\mathcal{E}|\mathcal{C}_0$\texttt{]}, so as to reduce (in b) and c)) to the case where $\operatorname{card}\mathcal{C} \leqslant \pi$.

\textbf{Corollary 4.14.19.} \textit{Under the conditions of 4.14.17, the subcategory $\uGamma_{\cocart}(\mathcal{E}/\mathcal{C})$ of $\mathcal{M} = \uGamma(\mathcal{E}/\mathcal{C})$ (case $p$ cofibrant), resp. the subcategory $\uGamma_{\cart}(\mathcal{E}/\mathcal{C})$ (case $p$ fibrant) is accessible in $\mathcal{M}$, hence accessible and the inclusion functor} \texttt{[is]} \textit{accessible.}

(\textbf{N.B.} These two categories are strictly full in $\uGamma(\mathcal{E}/\mathcal{C})$, so the assertion that they are accessible in it \textit{is equivalent} to the last assertion.)

\textsc{Proof of} 4.14.19. We may assume $\mathcal{C}$ reduced, hence small. Let us consider the case where $p$ is cofibrant. For every arrow $\lambda : s \longrightarrow t$ in $\mathcal{C}$, let us consider the functor
\[
\varphi_\lambda \ : \ \mathcal{M} \longrightarrow \uFl(\mathcal{E}_t)
\]
defined by
\msp{page 394}
\[
\varphi_\lambda(f) \ = \ (\lambda_*(f(s)) \xrightarrow{f(\lambda)^\natural} f(t)) \ .
\]
It is an accessible functor, since its composite with the source and target functors are the functors $\lambda_*\varepsilon_s$ and $\varepsilon_t$, hence accessible, $\varepsilon_s$, $\varepsilon_t$ and $\lambda_*$ being accessible. Let us consider the strictly full subcategory
\[
\uIs(\mathcal{E}_t) \ \subset \ \uFl(\mathcal{E}_t) \ ,
\]
we know that it is accessible in $\uFl(\mathcal{E}_t)$ (4.14.12). Thus its inverse image under $\varphi_\lambda$ is accessible in $\mathcal{M}$, and consequently so is
\[
\uGamma_{\cart}(\mathcal{E}/\mathcal{C}) \ \overset{\text{def}}{=} \ \bigcap_{\lambda\in\Ob\mathcal{C}} \varphi_\lambda^{-1}(\uIs(\mathcal{E}_t))
\]
% [sic] \uGamma_{\cart} -> \uGamma_{\cocart}, the case treated is p cofibrant
% [sic] \lambda\in\Ob\mathcal{C} -> \lambda\in\Fl\mathcal{C}, \lambda runs over arrows of \mathcal{C}
as the intersection of a \textit{small} family of subcategories accessible in $\mathcal{M}$, q.e.d. The case where $p$ is fibrant is treated in a similar way, by introducing the functors
\[
\psi_\lambda \ : \ \mathcal{M} = \uGamma(\mathcal{E}/\mathcal{C}) \longrightarrow \uFl(\mathcal{E}_s)
\]
given by
\[
\psi_\lambda(f) \ = \ (f(s) \xrightarrow{f(\lambda)^\flat} \lambda^* f(t)) \ .
\]

\msp{page 395}
\textbf{Corollary 4.14.20.} \textit{Let $\pi_0$ be an infinite cardinal chosen as in 4.14.18 a).}
\begin{itemize}[leftmargin=2.5em]
\item[\circled{a}] \textit{Then $\mathcal{M}_{\cc} = \uGamma_{\cocart}(\mathcal{E}/\mathcal{C})$ (resp. $\mathcal{M}_{\cc} = \uGamma_{\cart}$) is $L_{\pi_0}$-stable in $\mathcal{M} = \uGamma(\mathcal{E}/\mathcal{C})$, or what amounts to the same thing (as it is strictly full), it satisfies $L_{\pi_0}$, and the inclusion functor satisfies $L_{\pi_0}$. The $\varinjlim$ large relative to $\pi_0$ in this category are computed argument by argument.}
\item[\circled{b}] \textit{For every cardinal $\pi \geqslant \pi_0$, every object of $\mathcal{M}_{\cc}$ which is $\pi$-accessible in $\mathcal{M}$ is so (tautologically) in $\mathcal{M}_{\cc}$. If $\pi \geqslant \sup(\pi_0, \cardred\mathcal{C})$, then (by 4.14.18 b) every object $f$ of $\mathcal{M}_{\cc}$ such that $f(s) \in \Ob\mathcal{E}_{s\,\pi}$ $\forall\ s \in S$} \texttt{[rather} $s \in \Ob C$\texttt{]} \textit{is in $\mathcal{M}_\pi$} \textup{and} \textit{in $\mathcal{M}_{\cc\,\pi}$.}
% foot: Version 17 décembre 2025

\printedpage{150}
% head: A. Grothendieck : Dérivateurs, Ch. XVIII, Édition des Universités de Montpellier II et Paris VII

\item[\circled{c}] \textit{Suppose $\pi_0$ chosen as in 4.14.18 c, and in addition that $\mathcal{M}_{\cc}$ is $\pi_0$-accessible in $\mathcal{M}$. (Such a $\pi_0$ exists.) We thus have, for $\pi \geqslant \pi_1 = \sup(\pi_0, \cardred\mathcal{C}, \cardred\mathcal{M}_\pi)$ such that $\pi^{\pi_0} = \pi$,}
\[
(\mathcal{M}_{\cc})_\pi \ = \ \mathcal{M}_\pi{\cap}\mathcal{M}_{\cc} \ = \quad
\begin{array}{l}
\textit{Strictly full subcategory of $\mathcal{M}_{\cc}$ consisting of the $f$}\\
\textit{such that $f(s) \in \Ob\mathcal{E}_{s\,\pi}$ $\forall\ s \in \Ob\mathcal{C}$.}
\end{array}
\]

\msp{page 396}
\textit{If in addition $\pi$ is chosen so that we have $\pi \geqslant \sup_{s\in\Ob\mathcal{C}} \cardred\mathcal{E}_{s\,\pi_0}$, then}
\[
\left\{
\begin{array}{rrcl}
 & \mathcal{M}_{\cc\,\pi} & \xleftarrow{\ \sim\ } & \uGamma_{\cocart}(\mathcal{E}(\pi)/\mathcal{C})\\[1ex]
\textit{resp.} & \mathcal{M}_{\cc\,\pi} & \xleftarrow{\ \sim\ } & \uGamma_{\cart}(\mathcal{E}(\pi)/\mathcal{C}) \ ,
\end{array}
\right.
\]
\textit{hence ($\pi$ being useful for $\mathcal{M}_{\cc}$) an equivalence of categories}
\[
\begin{array}{rrcl}
 & \uGamma_{\cocart}(\mathcal{E}/\mathcal{C}) & \xleftarrow{\ \sim\ } & \mathrm{Ind}_\pi(\uGamma_{\cocart}(\mathcal{E}(\pi)/\mathcal{C}))\\
\textit{resp.}\quad & \uGamma_{\cart}(\mathcal{E}/\mathcal{C}) & \xleftarrow{\ \sim\ } & \mathrm{Ind}_\pi(\uGamma_{\cart}(\mathcal{E}(\pi)/\mathcal{C})).
\end{array}
\]
\end{itemize}

\textsc{Proof}. Everything is clear in the order in which it is given using the two preceding corollaries.

\textbf{Corollary 4.14.21.} \textit{Let us consider a diagram of functors}
\[
\begin{tikzcd}[column sep=small, row sep=small]
\mathcal{E}_1 \arrow[dr, "p_1"] & & \mathcal{E}_2 \arrow[dl, "p_2"'] \\
& \mathcal{E}_0 \ , &
\end{tikzcd}
\tag*{$(*)$}
\]
\textit{the categories and the functors being accessible.}

\msp{page 397}
\begin{itemize}[leftmargin=2.5em]
\item[\textit{a)}] \textit{The 2-fiber product}
\[
\mathcal{M}_0 \ = \ \mathcal{E}_1 \overset{2}{\times}_{\mathcal{E}_0} \mathcal{E}_2
\]
\textit{is accessible, and the canonical functors $\mathcal{M} \longrightarrow \mathcal{E}_1$, $\mathcal{M} \longrightarrow \mathcal{E}_2$ are accessible.} % [sic] \mathcal{M} \longrightarrow \mathcal{E}_1, \mathcal{M} \longrightarrow \mathcal{E}_2 -> \mathcal{M}_0 \longrightarrow \mathcal{E}_1, \mathcal{M}_0 \longrightarrow \mathcal{E}_2, assertion a) concerns \mathcal{M}_0
\item[\textit{b)}] (\footnote[169]{This is contained in 4.14.22.1 above (p. 406), applied to $\Pi = \mathcal{E}_1{\times}\mathcal{E}_2$, $\Pi' = \mathcal{E}_0$, and $F, G : \Pi \rightrightarrows \Pi'$ given by $F = p_1\circ\pr_1$ and $G = p_2\circ\pr_2$. On the other hand, a) and c) follow easily from b).}). \textit{The category $\mathcal{M}$ of triples $(x_1,x_2,u)$ with $x_1 \in \Ob\mathcal{E}_1$, $x_2 \in \Ob\mathcal{E}_2$, and}
\[
u \ : \ p_1(x_1) \ \longrightarrow \ p_2(x_2) \ ,
\]
\textit{is accessible, and the canonical functors $\mathcal{M} \longrightarrow \mathcal{E}_1$, $\mathcal{M} \longrightarrow \mathcal{E}_2$ are accessible.}
\item[\textit{c)}] \textit{The strictly full subcategory $\mathcal{M}_0$ of $\mathcal{M}$ is accessible in $\mathcal{M}$.}
\end{itemize}

\noindent\textsc{Proof}. Let $\mathcal{C}$ be the category
\[
\begin{tikzcd}[column sep=small]
1 \arrow[dr] & & 2 \arrow[dl] \\
& 0 \ , &
\end{tikzcd}
\]
and $\mathcal{E}$ the cofibered category over $\mathcal{C}$ defined by the diagram $(*)$. We are under the conditions of 4.14.17. The category $\Gamma = \uGamma(\mathcal{E}/\mathcal{C})$ is interpreted
% foot: Version 17 décembre 2025
\printedpage{151}
% head: A. Grothendieck : Dérivateurs, Ch. XVIII, Édition des Universités de Montpellier II et Paris VII
\msp{page 398}
as the category of quintuples $(x_0,x_1,x_2,u_1,u_2)$, with $x_i \in \Ob\mathcal{E}$ for $i \in \{0,1,2\}$, and $u_1, u_2 \in \Fl\mathcal{E}_0$,
\[
\begin{tikzcd}[column sep=small, row sep=small]
p_1(x_1) \arrow[dr, "u_1"] & & p_2(x_2) \arrow[dl, "u_2"'] \\
& x_0 \ . &
\end{tikzcd}
\tag*{$(**)$}
\]
By 4.14.18, this category $\uGamma$ is accessible. The category $\mathcal{M}$ is equivalent to the strictly full subcategory $\Gamma_0$ of $\Gamma$, consisting of the quintuples for which we have $u_2 \in \uIs\mathcal{E}_0$. We have indeed functors quasi-inverse to each other % [sic] \Gamma_0 -> \Gamma_1, the subcategory is called \Gamma_1 in what follows
\[
\begin{tikzcd}
\mathcal{M} \arrow[r, shift left, "\alpha"] & \Gamma_1 \arrow[l, shift left, "\beta"]
\end{tikzcd}
\]
given by
\[
\begin{aligned}
\alpha(x_1,x_2,u) \ &= \ (p_2(x_2),x_1,x_2,u,\id_{p_2(x_2)})\\
\beta(x_0,x_1,x_2,u_1,u_2) \ &= \ (x_1,x_2,u_2^{-1}u_1) \ ,
\end{aligned}
\]
we then have
\[
\begin{aligned}
\beta\alpha \quad &= \quad \id_{\mathcal{M}}\\
\alpha\beta \quad &\xrightarrow{\ \sim\ } \quad \id_{\Gamma_1} \ .
\end{aligned}
\]
On the other hand, $\Gamma_1$ is obviously accessible in $\mathcal{M}$, as the inverse image of $\uIs(\mathcal{E}_0)$ (which % [sic] accessible dans \mathcal{M} -> accessible dans \Gamma, \Gamma_1 is a subcategory of \Gamma
\msp{page 399}
is accessible in $\uFl(\mathcal{E}_0)$ under the accessible functor
\[
\uGamma \ \longrightarrow \ \uFl(\mathcal{E}_0)
\]
given by $(x_1, x_1, x_2, u_1, u_2) \longmapsto u_2$. This proves that $\mathcal{M}$, which is equivalent to $\Gamma_1$, is also accessible, and the accessibility of the canonical functor $\mathcal{M} \longrightarrow \mathcal{E}_i$ ($i \in \{1,2\}$) follows from that of $\Gamma \longrightarrow \mathcal{E}_i$, implying that of the composite $\Gamma_1 \hookrightarrow \Gamma \longrightarrow \mathcal{E}_i$. This therefore proves b). To prove a), it suffices to see that $\mathcal{M}_0$ is accessible in $\mathcal{M}$, now it is again the inverse image of $\uIs(\mathcal{E}_0) \subset \uFl(\mathcal{E}_0)$ under the canonical accessible functor % [sic] (x_1, x_1, x_2, u_1, u_2) -> (x_0, x_1, x_2, u_1, u_2), the quintuplets begin with x_0
\[
\mathcal{M} \ \longrightarrow \ \uFl(\mathcal{E}_0) \ ,
\]
so $\mathcal{M}$ is indeed accessible in $\mathcal{M}$, as the inverse image of a subcategory accessible in $\uFl(\mathcal{E})$. We can moreover also interpret $\mathcal{M}$, up to equivalence of categories, as % [sic] donc \mathcal{M} est bien accessible -> donc \mathcal{M}_0 est bien accessible, \mathcal{M}_0 is the subcategory in question
% [sic] \uFl(\mathcal{E}) -> \uFl(\mathcal{E}_0), as in the preceding display
% [sic] interpréter \mathcal{M} -> interpréter \mathcal{M}_0, as in the following display
\msp{page 400}
\[
\mathcal{M}_0 \ \simeq \ \uGamma_{\cocart}(\mathcal{E}/\mathcal{C}) \ ,
\]
the canonical functors $\mathcal{M}_0 \longrightarrow \mathcal{E}_i$ ($i \in \{0,1\}$) being interpreted as evaluation functors, which gives a) directly, as a special case of 4.14.19. % [sic] i \in \{0,1\} -> i \in \{1,2\}, the canonical functors of a) go to \mathcal{E}_1, \mathcal{E}_2
% foot: Version 17 décembre 2025

\printedpage{152}
% head: A. Grothendieck : Dérivateurs, Ch. XVIII, Édition des Universités de Montpellier II et Paris VII

\textbf{4.14.22.} \textsc{Proof of} 4.14.17. The case where $p$ is cofibrant has already been proved (th. 4.10.8, p. 263). We are going to see that the case where $p$ is fibrant reduces to the preceding theorem (in a very special case), using to the full the general results on accessible parts of an accessible category. Let us note that assertions a) and b) of corollary 4.14.18 are proved in AQT \texttt{[âne qui trotte: at a donkey's trot, i.e. by routine verification]} fashion without having to appeal to 4.14.17. Part a) already implies that the evaluation functors $\varepsilon_s : \mathcal{M} \longrightarrow \mathcal{E}_s$ are accessible. It therefore remains to prove
\msp{page 401}
that $\mathcal{M} = \uGamma(\mathcal{E}/\mathcal{C})$ is accessible, in the case where $p$ is \textit{fibrant}. We shall moreover assume that $\mathcal{C}$ is small (replacing $\mathcal{C}$ by an equivalent small full subcategory if need be).

Let
\[
\begin{aligned}
\Pi \ &= \ \prod\nolimits_{s\in\Ob\mathcal{C}} \mathcal{E}_s\\
\Pi' \ &= \ \prod\nolimits_{\lambda\in\Fl\mathcal{C}} \mathcal{E}_{s(\lambda)} \ ,
\end{aligned}
\]
where $s(\lambda)$ denotes the \textit{source} $s$ of the arrow $\lambda : s \longrightarrow t$ in $\mathcal{C}$. $\Pi$, $\Pi'$ are accessible categories, as (small) products of accessible categories. (Trivial special case of 4.10.8, for a discrete index category.) We are going to define two functors
\[
\Pi \ \overset{F}{\underset{G}{\rightrightarrows}} \ \Pi'
\]
by defining for every $\lambda : s_0 \longrightarrow t_0$ in $\Fl\mathcal{C}$ two functors
\msp{page 402}
\[
\left\{
\begin{array}{lcll}
\Pi & \overset{F_\lambda}{\underset{G_\lambda}{\rightrightarrows}} & \mathcal{E}_{s(\lambda)=s_0} & \\
F_\lambda & = & \pr_{s_0} & \text{projection of } \Pi = \prod\nolimits_s \mathcal{E}_s \text{ onto the factor } \mathcal{E}_{s_0}\\
G_\lambda & = & \lambda^*\pr_{s_1} & \text{composite } \Pi \xrightarrow{\pr_{s_1}} \mathcal{E}_{s_1} \xrightarrow{\lambda^*} \mathcal{E}_{s_0}\\
 & & & \text{\texttt{[rather} } \lambda^*\pr_{t_0} \text{ \texttt{resp.} } \Pi \xrightarrow{\pr_{t_0}} \mathcal{E}_{t_0} \xrightarrow{\lambda^*} \mathcal{E}_{s_0}\text{\texttt{]}}
\end{array}
\right.
\qquad \bigl(\mbox{\footnotemark[170]}, \mbox{\footnotemark[171]}\bigr) \ .
\]
\footnotetext[170]{It is immediate that the $F_\lambda$, $G_\lambda$, hence $F$ and $G$, are \textit{accessible} functors.}
\footnotetext[171]{\textbf{N.B.} We assume chosen for every $\lambda \in \Fl\mathcal{C}$ an inverse image functor $\lambda^* : \mathcal{E}_t \longrightarrow \mathcal{E}_s$. We assume that if $\lambda$ is an identity arrow, $\id_s$, \texttt{[we]} choose $\lambda^* = \id_{\mathcal{E}_s}$.}

Let $f \in \mathcal{M} = \uGamma(\mathcal{E}/\mathcal{C})$, then for every $\lambda \in \Fl\mathcal{C}$, $f$ defines an arrow
\[
\alpha_\lambda(f) \ : \quad \underbrace{f(s)}_{= F_\lambda(\varepsilon(f))} \quad \xrightarrow{f(\lambda)^\flat} \quad \underbrace{\lambda^*(f(t))}_{= G_\lambda(\varepsilon(f))} \ ,
\]
\textit{functorial} in $f$. We have a functor
\[
\varepsilon \ = \ (\varepsilon_s)_{s\in\Ob\mathcal{C}} \ : \ \mathcal{M} \ \longrightarrow \ \Pi \ ,
\]
and for every $\lambda \in \Fl\mathcal{C}$ a functorial homomorphism
\[
\alpha_\lambda \ : \ F_\lambda \circ \varepsilon \ \longrightarrow \ G_\lambda \circ \varepsilon \ ,
\]
whence a functorial homomorphism
\[
\alpha \ = \ (\alpha_\lambda)_{\lambda\in\Fl\mathcal{C}} \ : \ F \circ \varepsilon \ \longrightarrow \ G \circ \varepsilon \ .
\]
Let then $\mathcal{M}'$ be the following category:
% foot: Version 17 décembre 2025
\printedpage{153}
% head: A. Grothendieck : Dérivateurs, Ch. XVIII, Édition des Universités de Montpellier II et Paris VII
\msp{page 403}
\[
\mathcal{M}' \quad = \quad \parbox[t]{0.5\textwidth}{category of pairs $(X,u)$, with $X$ in $\Ob\Pi$, and $u : F(X) \longrightarrow G(X)$ an arrow in $\Pi'$.}
\]
The functorial homomorphism $\alpha : F \circ \varepsilon \longrightarrow G \circ \varepsilon$ defines a functor
\[
\mathcal{M} \ \longrightarrow \ \mathcal{M}' \ ,
\]
and it is immediate that this functor is fully faithful. In fact, the category $\mathcal{M}'$ is isomorphic to the category consisting of the pairs $(f_0,f_1)$ of two maps
\[
\begin{array}{lcccccc}
f_0 : & \Ob\mathcal{C} & \longrightarrow & \Ob\mathcal{E} & s & \longmapsto & f_0(s)\\[1ex]
f_1 : & \Fl\mathcal{C} & \longrightarrow & \Fl\mathcal{E} & \lambda & \longmapsto & f_1(\lambda)
\end{array}
\]
in such a way that $f_1$ and $f_0$ are compatible with the source and target maps, i.e. that they define a homomorphism of \textit{diagrams} $f : \diag\mathcal{C} \longrightarrow \diag\mathcal{E}$, and that this homomorphism is a right inverse of $\mathrm{diag}(p) : \diag\mathcal{E} \longrightarrow \diag\mathcal{C}$.

\msp{page 404}
The preceding functor $\mathcal{M} \hookrightarrow \mathcal{M}'$ is interpreted as the functor which associates to every sections $f$ of $\mathcal{E}$ over $\mathcal{C}$, the corresponding morphism of diagrams. It is therefore the inclusion of a \textit{strictly full} subcategory of $\mathcal{M}'$. Thus to prove that $\mathcal{M}$ is accessible, it suffices to prove: % [sic] toute sections -> toute section, agreement of toute
\begin{itemize}[leftmargin=2.5em]
\item[a)] $\mathcal{M}'$ is accessible.
\item[b)] $\mathcal{M}$ is accessible in $\mathcal{M}'$.
\end{itemize}
Let us first prove b) (admitting a). The subcategory $\mathcal{M}$ of $\mathcal{M}'$ is defined by two types of conditions, on an object $f = (f_0,f_1)$ of $\mathcal{M}'$:
\begin{itemize}[leftmargin=2.5em]
\item[1°)] $\forall\ s \in \Ob\mathcal{C}$, $f_1(\id_s) = \id_{f_0(s)}$.
\item[2°)] If $s \xrightarrow{\lambda} t \xrightarrow{\mu} u$ are two composable arrows in $\mathcal{C}$, then
\[
f_1(\mu\lambda) \ = \ f_1(\mu)f_1(\lambda) \ .
\]
\end{itemize}
These conditions are made explicit the one and \texttt{[the other by]} the equality of a double arrow in $\mathcal{E}_s$,
\msp{page 405}
namely
\[
f_0(s) \overset{f_1(\id_s)}{\underset{\id_{f_0(s)}}{\rightrightarrows}} f_0(s)
\]
in case 1°), and
\[
f_0(s) \overset{f_1(\mu\lambda)^\flat}{\underset{(f_1(\mu)f_1(\lambda))^\flat}{\rightrightarrows}} (\mu\lambda)^*(f_0(u))
\]
in case 2°). In either case, we have a functor
\[
\mathcal{M}' \ \xrightarrow{\gamma_s \text{ resp. } \gamma_{\lambda,\mu}} \ \uHom(I,\mathcal{E}_s)
\]
% foot: Version 17 décembre 2025

\end{document}
